\chapter{Resultados}

Este capítulo apresenta os resultados obtidos na análise empírica, iniciando pela caracterização 
das propriedades estatísticas da série de retornos cambiais e de sua volatilidade. Em seguida, 
são apresentados e discutidos os resultados das previsões realizadas pelos modelos Random Forest, 
GARCH e EGARCH, considerando as métricas de desempenho e os testes de comparação preditiva. Por fim, 
é realizada uma análise de robustez por meio da substituição da janela expansiva por uma janela rolante 
de cinco anos, com o objetivo de avaliar a sensibilidade dos resultados a essa alteração na estratégia de estimação.

\section{Análise descritiva}

Esta seção apresenta a análise descritiva da série de retornos log-diários da taxa de câmbio
USD/BRL, cobrindo o período de janeiro de 2000 a maio de 2026. O objetivo é caracterizar
empiricamente as propriedades estatísticas da série e validar os chamados fatos estilizados das
finanças, que fundamentam a escolha de modelos adotados.

O gráfico da série de câmbio USD/BRL em nível mostra uma trajetória ascendente ao longo do
período amostral, como mostra a figura 5.1:

\begin{figure}[!htbp]
    \centering
    \includegraphics[
    width=0.65\textwidth,
    height=0.4\textheight,
    keepaspectratio
]{../../../OUTPUTS/figures/serie_cambio.png}
    \caption{Série do càmbio}
    \label{fig:previsoes}
\end{figure}

Partindo de valores próximos a \text{R\$ 1,80/US\$} no início do ano de 2000, a taxa atingiu níveis
superiores a \text{R\$ 6,00/US\$} nos anos recentes, evidenciando uma tendência estocástica de longo
prazo. A decomposição sazonal da série, apresentada na Figura A.1 do Apêndice A, confirma visualmente a
presença de tendência crescente e componente sazonal, com o resíduo concentrando os
movimentos abruptos associados a eventos de estresse. Esses elementos são consistentes com a
hipótese de que a série de câmbio é não estacionária em nível, o que é verificado formalmente pelo
teste ADF na tabela 5.1.


\begin{figure}[htbp]
    \centering
    \includegraphics[width=0.9\textwidth]{../../../OUTPUTS/figures/volatilidade_movel_21d.png}
    \caption{Volatilidade móvel de 21 dias}
    \label{fig:previsoes}
\end{figure}

A série de retornos log-diários (Figura A.3 do Apêndice A) apresenta comportamento distinto do nível da taxa de câmbio, 
com oscilações em torno de zero e variância que se altera ao longo do tempo. Esse comportamento também pode ser observado 
na volatilidade histórica móvel de 21 dias (Figura 5.2), que evidencia episódios de maior volatilidade, especialmente em 2002, 
2008, 2015–2016 e 2020, intercalados por períodos de relativa calmaria. Essa alternância, denominada clustering de volatilidade, 
é um dos fatos estilizados das séries financeiras e constitui uma das principais motivações para a adoção de modelos da família GARCH.

\begin{table}[htbp]
    \centering
    \caption{Testes de estacionariedade e heterocedasticidade}
    \label{tab:testes_principais}
    \begin{tabular}{lrr}
        \toprule
        Teste & Estatística & $p$-valor \\
        \midrule
        ADF (câmbio em nível) & -0,760 & 0,8305 \\
        ADF (retornos) & -14,025 & $<0,001$ \\
        Ljung--Box (retornos ao quadrado, lag 10) & 4685,152 & $<0,001$ \\
        ARCH & 1669,466 & $<0,001$ \\
        \bottomrule
    \end{tabular}

    \vspace{0.2cm}
    \begin{minipage}{0.85\textwidth}
        \footnotesize
        \textit{Nota:} O teste ADF tem como hipótese nula a presença de raiz unitária. O teste Ljung--Box testa a ausência de autocorrelação até a defasagem considerada, enquanto o teste ARCH testa a ausência de efeitos ARCH.
    \end{minipage}
\end{table}

Os testes de estacionariedade indicam que não se rejeita a hipótese nula de raiz unitária para a taxa de câmbio em nível (p-valor = 0,831), 
enquanto essa hipótese é rejeitada para os retornos log-diários (p-valor < 0,001). Esses resultados são compatíveis com a não estacionariedade 
da série em nível e a estacionariedade dos retornos, que são utilizados nas etapas seguintes da análise.

Conforme apresentado na Tabela A.2 do Apêndice A, a distribuição dos retornostem média próxima de zero (0,000156) e baixa assimetria (0,026), mas curtose de 7,994, indicando uma distribuição 
com caudas mais pesadas diferentes de uma distribuição normal. O teste de Jarque-Bera rejeita a hipótese de normalidade (p-valor < 0,001), resultado 
também ilustrado pela comparação entre a densidade empírica e a normal teórica na Figura A.5 do Apêndice A. Além disso, o teste t não rejeita a hipótese 
de média igual a zero (p-valor = 0,205). Em conjunto, esses resultados mostram como funciona a distribuição dos retornos e são relevantes para a especificação dos modelos de volatilidade.

Por fim, o teste de Ljung-Box aplicado aos retornos ao quadrado rejeita a hipótese de ausência de autocorrelação tanto para 10 quanto para 20 defasagens 
(p-valores < 0,001). O teste ARCH também rejeita a hipótese nula de ausência de efeitos ARCH (p-valor < 0,001). Essas evidências apontam para dependência temporal na volatilidade 
e sustentam o uso de modelos da família GARCH. Os correlogramas dos retornos e dos retornos ao quadrado, apresentados na Figura A.2 do Apêndice A, complementam essa análise. 
Os resultados completos dos testes estatísticos são apresentados na Tabela A.1 do Apêndice A.

\section{Resultados principais}

\begin{table}[htbp]
    \centering
    \caption{Desempenho preditivo dos modelos}
    \label{tab:metricas}
    \begin{tabular}{lrrrr}
        \toprule
        Modelo & MSE & RMSE & MAE & QLIKE \\
        \midrule
        RF     & $2{,}604 \times 10^{-8}$ & 0,000161 & 0,000093 & $-8{,}457$ \\
        EGARCH & $2{,}643 \times 10^{-8}$ & 0,000163 & 0,000092 & $-8{,}451$ \\
        GARCH  & $2{,}664 \times 10^{-8}$ & 0,000163 & 0,000094 & $-8{,}448$ \\
        \bottomrule
    \end{tabular}

    \vspace{0.2cm}
    \footnotesize Fonte: Elaboração própria.
\end{table}

A Tabela 5.2 apresenta as métricas de desempenho preditivo dos modelos estimados. De maneira geral, o Random Forest apresentou os melhores resultados em termos de MSE, 
RMSE e QLIKE, enquanto o EGARCH obteve o menor MAE. Embora as diferenças entre os valores sejam relativamente pequenas, o desempenho do Random Forest indica uma vantagem 
numérica na previsão da variância cambial, especialmente considerando as métricas que penalizam erros de previsão de forma mais sensível como o MSE.

Em relação aos modelos da família GARCH, o EGARCH apresentou resultados ligeiramente superiores ao GARCH nas quatro métricas consideradas, com exceção do MAE, no qual 
também obteve o melhor resultado entre os três modelos. Esse resultado sugere que a especificação assimétrica do EGARCH pode oferecer alguma vantagem na representação da 
dinâmica da volatilidade cambial, embora essa diferença seja pequena.

\begin{figure}[!htbp]
    \centering
    \includegraphics[
        width=0.9\textwidth,
        keepaspectratio
    ]{../../../OUTPUTS/figures/actual_vs_pred_garch_expansiva.png}
    \caption{Variância observada e prevista pelo modelo GARCH(1,1) — janela expansiva}
    \label{fig:previsoes_garch}
\end{figure}

\begin{figure}[!htbp]
    \centering
    \includegraphics[
        width=0.9\textwidth,
        keepaspectratio
    ]{../../../OUTPUTS/figures/actual_vs_pred_egarch_expansiva.png}
    \caption{Variância observada e prevista pelo modelo EGARCH(1,1) — janela expansiva}
    \label{fig:previsoes_garch}
\end{figure}

\begin{figure}[!htbp]
    \centering
    \includegraphics[
        width=0.9\textwidth,
        keepaspectratio
    ]{../../../OUTPUTS/figures/actual_vs_pred_rf_expansiva.png}
    \caption{Variância observada e prevista pelo modelo Random Forest — janela expansiva}
    \label{fig:previsoes_garch}
\end{figure}

A análise gráfica das variâncias observadas e previstas, apresentada nas Figuras 5.3 a 5.5, permite complementar a comparação. É possível notar que os três 
modelos conseguem acompanhar parcialmente a dinâmica da volatilidade ao longo do período analisado, identificando momentos de maior instabilidade. Entretanto, 
é possível identificar, também, que o Random Forest apresenta uma capacidade relativamente maior de acompanhar alguns dos picos de volatilidade, especialmente 
nos períodos de maior instabilidade. Apesar disso, as previsões dos três modelos apresentam um comportamento mais suavizado em comparação com a variância observada, 
evidenciando dificuldades em representar integralmente os movimentos extremos. Assim, embora o Random Forest apresente uma resposta visualmente mais próxima da variância 
observada em determinados momentos, essa característica não se traduz em uma superioridade estatisticamente significativa, conforme indicado pelos testes de Diebold-Mariano.

Além das métricas de desempenho, a tabela B.3 no Apêndice B mostra a importância das variáveis utilizadas pelo modelo Random Forest, com o objetivo de identificar quais atributos contribuíram 
mais para suas previsões. Na janela expansiva, a média móvel dos retornos absolutos em cinco dias apresentou a maior importância (0,163), seguida pelo retorno defasado em um 
período (0,118) e pela volatilidade defasada em um período (0,105). Também se destacaram a média móvel dos retornos absolutos em 21 dias e a volatilidade histórica de 21 dias. 
Esses resultados indicam a relevância das informações recentes de retorno e das medidas de volatilidade passada para as previsões do modelo.

\begin{table}[!htbp]
    \centering
    \caption{Teste Diebold--Mariano para a função de perda MSE}
    \label{tab:dm_mse}
    \begin{tabular}{lrrr}
        \toprule
        Modelo A & Modelo B & DM-HLN & $p$-valor \\
        \midrule
        RF    & GARCH  & -1,128 & 0,259 \\
        RF    & EGARCH & -0,700 & 0,484 \\
        GARCH & EGARCH &  0,932 & 0,351 \\
        \bottomrule
    \end{tabular}
    \begin{center}
        \footnotesize Fonte: Elaboração própria.
    \end{center}
\end{table}

\begin{table}[!htbp]
    \centering
    \caption{Teste Diebold--Mariano para a função de perda QLIKE}
    \label{tab:dm_qlike}
    \begin{tabular}{lrrr}
        \toprule
        Modelo A & Modelo B & DM-HLN & $p$-valor \\
        \midrule
        RF    & GARCH  & -0,827 & 0,408 \\
        RF    & EGARCH & -0,525 & 0,600 \\
        GARCH & EGARCH &  0,440 & 0,660 \\
        \bottomrule
    \end{tabular}
    \begin{center}
        \footnotesize Fonte: Elaboração própria.
    \end{center}
\end{table}

Por fim, os resultados dos testes de Diebold-Mariano com a correção de Harvey, Leybourne e Newbold, apresentados nas Tabelas 5.3 e 5.4, não apontaram diferenças 
estatisticamente significativas entre os modelos, tanto para a função de perda MSE quanto para a QLIKE. Em todas as comparações, os p-valores foram superiores a 0,05, 
ou seja, não foi possível rejeitar a hipótese nula de igualdade de desempenho preditivo. Dessa forma, apesar da vantagem numérica observada para o Random Forest em 
algumas métricas, não há evidências estatísticas suficientes para afirmar que seu desempenho preditivo seja superior ao dos modelos GARCH e EGARCH.

Em conjunto, os resultados indicam que tanto os modelos econométricos da família GARCH quanto o Random Forest foram capazes de produzir previsões competitivas para 
a volatilidade cambial. A ausência de diferenças estatisticamente significativas sugere que a maior flexibilidade do Random Forest não se traduziu, nessa análise, 
em uma superioridade preditiva estatisticamente comprovada.

Resultados complementares das estimações dos modelos GARCH e EGARCH, incluindo os coeficientes estimados e suas respectivas estatísticas, são apresentados nas 
Seções B.1 e B.2 do Apêndice B. O apêndice também apresenta a importância das variáveis utilizadas pelo Random Forest (Seção B.3), permitindo identificar aquelas 
que mais contribuíram para as previsões, bem como os resultados completos dos testes Diebold-Mariano para as funções de perda MSE e QLIKE (Seção B.4).


\FloatBarrier
\section{Testes de robustez}

Para avaliar a robustez dos resultados obtidos na análise principal, foi realizada uma análise complementar utilizando uma janela rolante de cinco anos, em 
substituição à janela expansiva. Essa especificação permite verificar se o desempenho relativo dos modelos se mantém quando as estimações consideram apenas 
as observações mais recentes, atribuindo maior relevância às condições de mercado próximas ao período de previsão. A comparação foi realizada entre os modelos 
Random Forest e GARCH, conforme os resultados apresentados a seguir.

Na análise de robustez com janela rolante de cinco anos, o modelo EGARCH foi excluído devido à instabilidade observada no seu processo de estimação. Foram 
identificados episódios de não convergência do algoritmo de otimização e, mesmo em estimações com convergência reportada, ocorreram previsões de variância 
excessivamente elevadas. Esses problemas comprometeram a confiabilidade das previsões e motivaram a exclusão do modelo dessa especificação. Os diagnósticos 
das estimações são apresentados nas Tabelas B.6 e B.7 do Apêndice B.

\begin{table}[htbp]
    \centering
    \caption{Desempenho preditivo dos modelos — janela rolante}
    \label{tab:metricas_rolante}
    \begin{tabular}{lrrrr}
        \toprule
        Modelo & MSE & RMSE & MAE & QLIKE \\
        \midrule
        GARCH & $2{,}614 \times 10^{-8}$ & 0,000162 & 0,000091 & $-8{,}458$ \\
        RF    & $2{,}700 \times 10^{-8}$ & 0,000164 & 0,000095 & $-8{,}435$ \\
        \bottomrule
    \end{tabular}

    \vspace{0.2cm}
    \footnotesize Fonte: Elaboração própria.
\end{table}

A Tabela 5.5 apresenta as métricas de desempenho preditivo dos modelos estimados com janela rolante de cinco anos. De maneira geral, o GARCH apresentou 
os melhores resultados nas quatro métricas consideradas, obtendo menores valores de MSE, RMSE e MAE, além de um QLIKE mais negativo. Embora as diferenças entre 
os valores sejam relativamente pequenas, os resultados indicam uma vantagem numérica do GARCH em relação ao Random Forest na previsão da variância cambial.
Dessa forma, diferentemente do observado na janela expansiva, em que o Random Forest apresentou os melhores resultados em algumas métricas, a especificação 
com janela rolante favoreceu o modelo GARCH.

\begin{figure}[!htbp]
    \centering
    \includegraphics[
        width=0.9\textwidth,
        keepaspectratio
    ]{../../../OUTPUTS/figures/actual_vs_pred_garch_rolante.png}
    \caption{Variância observada e prevista pelo modelo GARCH(1,1) — janela rolante}
    \label{fig:previsoes_garch}
\end{figure}

\begin{figure}[!htbp]
    \centering
    \includegraphics[
        width=0.9\textwidth,
        keepaspectratio
    ]{../../../OUTPUTS/figures/actual_vs_pred_rf_rolante.png}
    \caption{Variância observada e prevista pelo modelo Random Forest — janela rolante}
    \label{fig:previsoes_garch}
\end{figure}

As Figuras 5.6 e 5.7 apresentam as variâncias observadas e previstas pelos modelos GARCH e Random Forest, respectivamente, considerando a janela rolante de 
cinco anos. De maneira geral, ambos os modelos conseguem acompanhar parte da dinâmica da volatilidade cambial, identificando períodos de maior instabilidade 
e apresentando previsões mais elevadas nos momentos de maior volatilidade observada.

Entretanto, é possível notar que as previsões dos dois modelos apresentam dificuldades em acompanhar a magnitude dos picos mais acentuados da variância observada. 
O GARCH apresenta uma trajetória de previsão relativamente mais suave, enquanto o Random Forest demonstra uma maior oscilação nas previsões, reagindo de maneira 
mais pronunciada a algumas variações da série. Apesar dessas diferenças visuais, ambos apresentam limitações na antecipação de movimentos extremos.

Além disso, a tabela B.8 do Apêndice B, mostra importância das variáveis utilizadas pelo modelo Random Forest, com o objetivo de identificar quais atributos 
contribuíram mais as previsões considerando a janela rolante. Conforme apresentado na tabela, os retornos defasados e as medidas de volatilidade passada tiveram 
destaque entre as dez variáveis mais importantes. Em particular, o retorno defasado em cinco períodos apresentou a maior importância (0,109), seguido pela volatilidade 
defasada em cinco períodos (0,100) e pelo retorno defasado em um período (0,100). Também se destacaram a volatilidade histórica de 21 dias e a média móvel dos retornos 
absolutos em 21 dias. Diferentemente da janela expansiva, em que as médias móveis dos retornos absolutos ocuparam posições de maior destaque, na janela rolante as variáveis 
de retorno e volatilidade defasados apresentaram maior importância relativa.

\begin{table}[!htbp]
    \centering
    \caption{Teste Diebold-Mariano para a janela rolante}
    \label{tab:dm_rolante}
    \begin{tabular}{lrrrr}
        \toprule
        Modelo A & Modelo B & Perda & DM-HLN & $p$-valor \\
        \midrule
        RF & GARCH & MSE   & 2,547 & 0,011 \\
        RF & GARCH & QLIKE & 2,376 & 0,018 \\
        \bottomrule
    \end{tabular}
    \begin{center}
        \footnotesize Nota: Ambas as comparações utilizam 1.602 observações.

        \footnotesize Fonte: Elaboração própria.
    \end{center}
\end{table}

Os resultados do teste Diebold-Mariano, apresentados na Tabela 5.6, indicam diferenças estatisticamente significativas entre os modelos para ambas as 
funções de perda. Para o MSE, a estatística DM-HLN foi de 2,547, com p-valor de 0,011, enquanto para o QLIKE foi de 2,376, com p-valor de 0,018. Em 
ambos os casos, a hipótese nula de igualdade de desempenho preditivo é rejeitada ao nível de significância de 5\%. Considerando que as diferenças médias 
de perda foram calculadas como RF menos GARCH, os resultados favorecem o modelo GARCH, corroborando sua superioridade observada nas métricas de desempenho.

Por fim, a análise de robustez mostra que a escolha da janela de estimação influencia o desempenho relativo dos modelos. Enquanto na especificação com janela 
expansiva o Random Forest apresentou os menores valores de MSE e QLIKE, embora sem diferenças estatisticamente significativas em relação aos demais modelos, 
na janela rolante de cinco anos o GARCH apresentou desempenho superior ao Random Forest nas duas funções de perda consideradas. Uma possível explicação é que 
a janela expansiva, ao incorporar um histórico mais amplo de observações, oferece ao Random Forest uma base maior para identificar padrões na série, enquanto 
a janela rolante restringe a estimação às informações mais recentes. Dessa forma, os resultados sugerem que a relação entre o desempenho preditivo e a estratégia 
de estimação merece atenção na escolha dos modelos de previsão da volatilidade cambial.